% META: {"id":"TCOMPL-INEG-EX-002","chapitre":"TCOMPL-INEGALITES","type_objet":"exercice","capacites":["TCOMPL-INEG-C3","TCOMPL-INEG-C5"],"capacites_codes":["C3","C5"],"methodes":[],"parcours":2,"competences":["calculer"],"duree_min":15,"mode_creation":"ex_nihilo","sources_inspiration":[],"fichier_tex":"chapitres/TCOMPL-INEGALITES/exercices/TCOMPL-INEG-EX-002.tex","corrige_tex":"chapitres/TCOMPL-INEGALITES/corriges/TCOMPL-INEG-CO-002.tex","status":"approved","origin":"generated","status_history":["generated","audited","accepted","approved"],"release_acceptance":"RELEASE_OWNER_BATCH_ACCEPTANCE","acceptance_closure_digest":"sha256:9b3ccf9a81c5520fbb7e03b7b2d3e7bf2057a02834b6d908bf3be5f49f3b553f","acceptance_receipt_digest":"sha256:4fad86f0282e0fa4e0419cca1ad6379ae7af5a280c04a4a90880f90a1c044242"}
% BEGIN-VERIFY
% from sympy import symbols, integrate, Rational
% x = symbols('x')
% L = (2*x**2 + x) / 3
% I = integrate(L, (x, 0, 1))
% assert I == Rational(7, 18)
% G = 1 - 2*I
% assert G == Rational(2, 9)
% END-VERIFY
\begin{exercice}{TCOMPL-INEG-EX-002}{2}{15}
On reprend le modèle $L(x) = \dfrac{2x^2+x}{3}$ de l'exercice précédent. Calculer l'indice de Gini associé à ce modèle.
\end{exercice}
